\subsubsection{Controle de atitude da espaçonave}
\label{sec:controle_PD_atitude}

A atitude da espaçonave é descrita por um quaternion unitário, conforme a Equação \eqref{eq:qrel_sincronizacao}, ou seja:

\begin{equation}
\vec q_B =
\begin{bmatrix}
q_{B,w} \\
q_{B,x} \\
q_{B,y} \\
q_{B,z}
\end{bmatrix}.
\end{equation}

Por sua vez, a referência de atitude é a atitude da espaçonave alvo, sendo especificada por um quaternion desejado, como:

\begin{equation}
\vec q_{B,\text{ref}} =
\begin{bmatrix}
q_{B,w,\text{ref}} \\
q_{B,x,\text{ref}} \\
q_{B,y,\text{ref}} \\
q_{B,z,\text{ref}}
\end{bmatrix}.
\end{equation}

O erro da atitude da espaçonave robótica é obtido através da Equação \eqref{eq:quat_relativo}, ou seja:

\begin{equation}
\label{eq:controle_PD_quat}
q_{rel} =
q_a \otimes q_p^{-1}.
\end{equation}

A lei de controle PD de atitude utilizada neste trabalho atua diretamente sobre a parte vetorial do quaternion de erro e sobre a velocidade angular:

\begin{equation}
\vec\tau_B=
-K_{p,B}
\begin{bmatrix}
q_{e,x} \\
q_{e,y} \\
q_{e,z}
\end{bmatrix}
-K_{d,B}
\begin{bmatrix}
\omega_x \\
\omega_y \\
\omega_z
\end{bmatrix},
\label{eq:controle_PD_atitude}
\end{equation}

em que \(K_{p,B}\) e \(K_{d,B}\) são matrizes de ganhos proporcionais e derivativos, respectivamente.
Por sua vez, a velocidade angular da base, expressa em $\{B\}$, é escrita como

\begin{equation}
\vec\omega_B =
\begin{bmatrix}
\omega_x \\
\omega_y \\
\omega_z
\end{bmatrix}.
\end{equation}

A ação de controle $\vec\tau_B$ é inserida no vetor de torques generalizados $\vec\tau$ das equações de movimento presentes nas Seções~\ref{sec:atitude_inercial} e~\ref{sec:dinAcoplada_equacoes_Lagrange}, especificamente nas Equações~\eqref{eq:euler_corpo_rigido} e~\eqref{eq:dinamica_geral}, já que as primeiras coordenadas generalizadas estão associadas à atitude da base, conforme a Seção~\ref{sec:var_estado_controle}.
Dessa forma, o controlador PD opera diretamente sobre a parcela correspondente aos torques da base no vetor $\vec\tau$.

Reorganizando os termos associados às coordenadas do \gls{mr} na Equação~\eqref{eq:dinAcoplada_base_compacta}, obtém-se a dinâmica da base explicitando o efeito de reação do manipulador:
\begin{equation}
M_{BB}(\vec q)\,\ddot{\vec q}_B
+
C_{BB}(\vec q,\dot{\vec q})\,\dot{\vec q}_B
=
\vec\tau_B
-
\underbrace{\bigl[
M_{B\theta}(\vec q)\,\ddot{\vec\theta}
+
C_{B\theta}(\vec q,\dot{\vec q})\,\dot{\vec\theta}
\bigr]}_{\displaystyle \vec\tau_{\text{reac}}},
\label{eq:dina_base_reac_compacta}
\end{equation}

em que

\begin{equation}
\vec\tau_{\text{reac}}
=
M_{B\theta}(\vec q)\,\ddot{\vec\theta}
+
C_{B\theta}(\vec q,\dot{\vec q})\,\dot{\vec\theta}
\label{eq:def_tau_reac}
\end{equation}

representa o torque de reação na base devido ao movimento das juntas do manipulador.
Durante o momento em que o \gls{mr} está fora de operação, ou seja, permanece travado em uma configuração fixa,

\begin{equation}
\dot{\vec\theta} = \vec 0,
\qquad
\ddot{\vec\theta} = \vec 0,
\end{equation}

tem-se $\vec\tau_{\text{reac}} = \vec 0$ e a Equação~\eqref{eq:dina_base_reac_compacta} reduz-se à dinâmica de um corpo rígido, em conformidade com a Equação~\eqref{eq:euler_corpo_rigido}.
Quando o \gls{mr} se movimenta, os termos de acoplamento $M_{BM}$ e $C_{BM}$ passam a introduzir $\vec\tau_{\text{reac}}$ como perturbação interna, de modo que o controlador PD em \eqref{eq:controle_PD_atitude} deve simultaneamente reduzir o erro de atitude e compensar os torques de reação induzidos pelo manipulador, preservando o alinhamento da base requerido para a operação de ancoramento (\emph{berthing}).